\section{Related Work}

\subsection{Classical Methods}

Classical methods apply a \emph{fixed transform} to the raw signal, and the dominant line has looked for slowly varying components on the premise that they reflect the factors that persist. STL \citep{cleveland1990} and EMD \citep{huang1998} split the waveform itself, and so deliver components of the signal rather than the factors behind it. Slow Feature Analysis instead optimizes slowness directly, minimizing the temporal derivative of its output under unit variance and pairwise decorrelation, with a polynomial expansion of the input making it nonlinear \citep{wiskott2002}; strengthening decorrelation into independence separates the slow components from one another as well \citep{blaschke2007}. These two are our label-free post-hoc rotation baselines. Predictable Feature Analysis keeps SFA's extraction scheme and swaps the criterion, selecting features best predicted from their own last $p$ values by a linear autoregressive model \citep{richthofer2013}. Reading the signal elsewhere reaches factors carried on the \emph{modulation} rather than the values, as in envelope analysis for rotating-machinery diagnosis \citep{randall2011}.

What the family shares is that \emph{the transform is not learned}: which statistics and bands to keep are fixed in advance, and the projection is refitted whenever the domain changes.

\subsection{Self-Supervised Learning for Time Series}

These methods differ chiefly in the \emph{form of the target}. Reconstruction is the default, and the prior a variational model places over the latents doubles as a handle on latent structure --- weighting the KL term pushes coordinates toward independence \citep{higgins2017}, a lineage leading to Section~\ref{sec:disent}; masked-patch pretraining is the current standard form \citep{nie2023}. Contrastive methods instead fix which pairs count as similar: CPC predicts the future contrastively in latent space \citep{oord2018}, and later work refines the pairs, contrasting hierarchically over augmented context views \citep{yue2022} or using local smoothness of the generative process to define neighborhoods with stationary properties \citep{tonekaboni2021}. What such a method discards is decided by the pair-selection rule and is not visible in the objective. Latent prediction takes the representation itself as the target \citep{lecun2022, assran2023}; its time-series variants differ in where that target is read --- masked patches \citep{ennadir2025}, the next latent token \citep{chemeris2026}, or an attention-pooled summary of the whole future interval \citep{petersen2026}.

This work sits in the latent-prediction family and carries both a regression ($L_1$) and a contrastive (InfoNCE) loss, to show the horizon gate does not depend on the loss form. In all three groups the representation is learned but the objective says nothing about what each coordinate means.

\subsection{Disentangled Representations for Time Series}
\label{sec:disent}

Success here has concentrated on factors \emph{constant over an entire sequence}, since invariance transfers directly into an objective: FHVAE assigns different priors to what is fixed within a sequence and what changes across segments \citep{hsu2017}, and the disentangled sequential autoencoder splits the latent into static and dynamic parts \citep{li2018}. On identifiability, learning to distinguish time segments exploits non-stationarity and, combined with linear ICA, identifies a nonlinear ICA model up to a pointwise transformation of each factor \citep{hyvarinen2016}; a sparse prior on the change between adjacent observations followed \citep{klindt2021}. The difficulty of model selection without labels \citep{locatello2019} applies to our hyperparameters too, and we return to it in the limitations.

What this leaves empty is the factor that varies slowly without being constant: the static/dynamic dichotomy files it under ``dynamic'' and stops. None of the three strands produces coordinate-block structure of the kind we want, and this work fills that place with the prediction horizon.
